\section{A conditional principal-norm rigidity route}\label{sec:rigidity}
\noindent\textit{Source anchor S: v0.10--v0.12. This is a theorem about a supplied smooth or holomorphic relative multiplier; the gamma comparison has not supplied such a multiplier with all its hypotheses.}

\subsection{The actual two commuting actions}
On $X=\mathbb C^3/\mathbb Z^3\simeq(\mathbb C^*)^3$ and its real torus $\mathbb T^3$, use
\[
 M=\begin{pmatrix}0&1&0\\-4&2&15\\-1&0&4\end{pmatrix},\quad
 N=2I-M=\begin{pmatrix}2&-1&0\\4&0&-15\\1&0&-2\end{pmatrix},\quad
 A=M^T,\quad B=N^T.
\]
The period row $(\tau,s,1)$ satisfies $(\tau,s,1)M=\epsilon(\tau,s,1)$ and $(\tau,s,1)N=\beta(\tau,s,1)$. We have $MN=NM$, $\det M=1$, $\det N=7$ and $N^3=7I$. Thus $N$ is a degree-seven cover on $X$. On the compact lattice $\mathbb Z[\omega]^3$ its degree is 49; these refer to different lattices.

The kernel on $X$ has representatives $j(1,2,4)^T/7$, $0\le j\le6$. The conductor point $p_3=e_2/3$ obeys $Mp_3-p_3=(0,5,1)^T$ and $Np_3-p_3=(0,-5,-1)^T$. The dual character $e_k(v)=e^{2\pi ik^Tv}$ satisfies $e_k(Mv)=e_{Ak}(v)$. Frequency transport is distinct from the scalar return sequence $F(\rho^m/3)$.

The cubic characteristic polynomial $(x-2)^3+7$ gives a stable real eigenvalue $a=2-7^{1/3}\in(0,1)$ and a conjugate unstable pair of modulus $a^{-1/2}>1$. Every nonzero integer $A$-orbit escapes in both directions: otherwise an integer vector would tend to zero along positive or negative iterates. Fourier orbit sums are therefore absolutely convergent for smooth coefficients. The dual degree-seven image is exactly
\[
 k\in B\mathbb Z^3\quad\Longleftrightarrow\quad
 \ell(k):=k_0+2k_1+4k_2\equiv0\pmod7.
\]
Indeed $k-\ell(k)e_0=B(-k_1,k_2,-8k_2)^T$ and $B(0,-2,15)^T=7e_0$. If $\ell(k)=7r$ its explicit preimage is $(-k_1,k_2-2r,15r-8k_2)^T$.

\subsection{The convergent transfer proof}
For functions on $\mathbb T^3$, let
\[
 S_Nh(v)=\sum_{\kappa\in\ker N}h(v+\kappa),\quad
 C_Nh(v)=h(Nv),\quad L_Nh(v)=\sum_{Nw=v}h(w),\quad
 \delta_Mh=h-h\circ M.
\]
The transfer $L_N$ is unnormalized. For $h=\sum c_ke_k$,
\[
 S_Nh=\sum_{k\in B\mathbb Z^3}7c_ke_k,\qquad
 L_Nh=\sum_k7c_{Bk}e_k,\qquad S_Nh=(L_Nh)\circ N,
 \qquad L_NC_Nh=7h.
\]
Since $M$ permutes $\ker N$ and commutes with $N$, the required operators commute with composition by $M$.

\begin{theorem}[Principal-norm rigidity]\label{thm:principal-rigidity}
Let $h,r$ be smooth mean-zero complex functions on $\mathbb T^3$. If
\begin{equation}\label{eq:norm-coboundary}
 S_Nh-h\circ N=\delta_Mr,
\end{equation}
then $h=\delta_Mg$ for a smooth mean-zero $g$, with convergent formula
\[
 g=-\frac17\sum_{j\ge1}L_N^jr.
\]
If $h,r$ extend holomorphically to all of $X$, so does $g$.
\end{theorem}
\begin{proof}
Apply $L_N/7$ to \eqref{eq:norm-coboundary}. The trace relation gives $L_Nh-h=(1/7)\delta_M L_Nr$, and induction yields
\[
 L_N^nh-h=\frac17\delta_M\sum_{j=1}^{n}L_N^jr.
\]
Because $B^3=7I$,
\[
 L_N^{3n}f=\sum_{k\ne0}343^nc_{7^nk}e_k
\]
for a mean-zero $f$. For any $C^q$ norm choose $s>\max(q,3)$. Smooth Fourier coefficients have finite weighted $\ell^1$ norms of every order, and the displayed expression is bounded, up to a derivative convention constant, by
\[
 7^{n(3-s)}\sum_{k\ne0}\|k\|^s|c_k|.
\]
It tends geometrically to zero and is summable. The other two residue classes are obtained by the fixed continuous operators $L_N,L_N^2$. The series for $g$ converges in every $C^q$ norm; taking the limit gives $h=\delta_Mg$. For a global holomorphic function on $X$, its Laurent coefficients decay faster than every fixed exponential. On each compact product annulus the decay at $7^nk$ dominates both $343^n$ and the annular factor $e^{R\|k\|}$, proving normal convergence there and a global holomorphic $g$.
\end{proof}
For a finite Laurent $r$ of coordinate height at most $H$, choose $n$ with $7^n>H$. Then $L_N^{3n}r=0$, so the series terminates by a derived bound. The finite coefficient recurrence has an integer Lean certificate; the smooth convergence and geometric interpretation above are ordinary analytic results.

\subsection{Phase-only compatibility and the missing identification}
Suppose a genuine nonzero relative multiplier has been constructed as
\[
 a_\Delta(v)=\alpha e^{2\pi im^Tv}e^{h(v)},\qquad
 m\in\mathbb Z^3,\quad h\text{ mean zero},
\]
and its based ratio is actually identified with $U_\Gamma/U_{\rm alg}$. Write $h^\#(v)=\overline{h(\overline v)}$. It suffices to establish \eqref{eq:norm-coboundary} for $h-h^\#$. The theorem gives $h-h^\#=g-g\circ M$. The gauge $e^{-g/2}$ replaces the periodic logarithm by $(h+h^\#)/2$, real on the real torus. Since $Mp_3=p_3$ modulo the lattice, this periodic gauge has equal values at the two endpoints, leaving
\[
 \frac{a_\Delta(p_3)}{a_\Delta(0)}=e^{2\pi im_2/3}e^{\text{real number}}.
\]
After the based identification the exact modulus equality forces that real exponent to vanish. The ratio then has cube one, and an independent branch enclosure can select one. The constant $\alpha$ cancels in this based ratio; it cannot simply be discarded in the unbased principal norm, where it contributes $\alpha^6$.

Neither the required periodic relative multiplier nor its phase-only norm relation is supplied by \eqref{eq:regularized-conjugation}. Scalar-plane divisor repair, periodic descent, based arithmetic identification and the norm equation are separate mathematical inputs. This conditional route is retained to locate the missing implication, not applied as if these inputs were proved.
